\documentclass[12pt,a4paper]{article}

\usepackage[a4paper,landscape,margin=1.4cm]{geometry}
\usepackage{fontspec}
\setmainfont{DejaVu Serif}
\usepackage{polyglossia}
\setdefaultlanguage{russian}
\usepackage{tikz}
\usetikzlibrary{arrows.meta}
\pagestyle{empty}

\begin{document}
\begin{center}
    {\LARGE\bfseries Архитектура решателя уравнений}
\end{center}
\vspace{-0.25cm}

\begin{center}
\begin{tikzpicture}[
    class/.style={
        draw=blue!45!black,
        fill=blue!3,
        rounded corners=2pt,
        line width=0.7pt,
        align=left,
        font=\footnotesize,
        text width=7.6cm,
        inner sep=8pt
    },
    abstract/.style={fill=blue!8},
    collection/.style={fill=teal!5},
    inheritance/.style={-{Triangle[open, length=3mm, width=3mm]}, line width=0.8pt},
    dependency/.style={dashed, -{Stealth[length=2mm]}, line width=0.8pt}
]
    \node[class] (complex) at (-10.7, 3.9) {
        \textbf{Complex (immutable)}\\[2pt]
        \rule{7.3cm}{0.4pt}\\
        - real: double\\
        - imag: double\\[2pt]
        \rule{7.3cm}{0.4pt}\\
        + operator+(rhs: Complex): Complex\\
        + operator-(rhs: Complex): Complex\\
        + operator*(rhs: Complex): Complex\\
        + operator/(rhs: Complex): Complex\\
        + operator==(rhs: Complex): bool\\
        + operator!=(rhs: Complex): bool
    };

    \node[class, abstract] (equation) at (0, 3.9) {
        \textbf{Equation} \textit{(abstract)}\\[2pt]
        \rule{7.3cm}{0.4pt}\\
        + operator()(x: Complex): Complex\\
        + solve(): vector<Complex>
    };

    \node[class] (linear) at (-4.7, -1.2) {
        \textbf{LinearEquation}\\[2pt]
        \rule{7.3cm}{0.4pt}\\
        - a: Complex\\
        - b: Complex\\[2pt]
        \rule{7.3cm}{0.4pt}\\
        + LinearEquation(a, b) {\(a \ne 0\)}\\
        + operator()(x: Complex): Complex\\
        + solve(): vector<Complex>
    };

    \node[class] (quadratic) at (4.7, -1.2) {
        \textbf{QuadraticEquation}\\[2pt]
        \rule{7.3cm}{0.4pt}\\
        - a: Complex\\
        - b: Complex\\
        - c: Complex\\[2pt]
        \rule{7.3cm}{0.4pt}\\
        + QuadraticEquation(a, b, c) {\(a \ne 0\)}\\
        + operator()(x: Complex): Complex\\
        + solve(): vector<Complex>
    };

    \node[class, collection] (collection) at (0, -6.0) {
        \textbf{EquationCollection}\\[2pt]
        \rule{7.3cm}{0.4pt}\\
        - equations: owned Equation objects\\[2pt]
        \rule{7.3cm}{0.4pt}\\
        + add(equation): void\\
        + solveAll()\\
        + evaluateAll(x: Complex)\\
        + findByRoot(x: Complex)
    };

    % Наследование: пустой треугольник указывает на базовый класс.
    \draw[inheritance] (linear.north east) -- (equation.south west);
    \draw[inheritance] (quadratic.north west) -- (equation.south east);

    % Equation использует Complex для аргумента и результата вычисления.
    \draw[dependency] (equation.west) -- node[above, fill=white, inner sep=1pt] {использует} (complex.east);

    % Коллекция владеет нулём или несколькими объектами Equation.
    \draw[{Diamond[length=3mm, width=3mm]}-] (collection.north) --
        node[pos=0.08, right, fill=white, inner sep=1pt] {1}
        node[pos=0.92, right, fill=white, inner sep=1pt] {0..*}
        (equation.south);
\end{tikzpicture}
\end{center}

\end{document}
